% =====================================================================
%  Section: Forced-response demonstration of the electrically driven
%           drivetrain torsional resonance (simplified + OpenFAST-FMU)
%  Required in preamble:
%    \usepackage{amsmath, amssymb}
%    \usepackage{graphicx}
%    \usepackage{booktabs}
%    \usepackage{siunitx}
%  Adjust the figure path (fig/...) to your own.
% =====================================================================

\section{Forced-Response Demonstration of the Torsional Resonance}
\label{sec:forcedresponse}

The modal analysis of \S\ref{sec:frequency-analysis} located a lightly damped
drivetrain torsional mode at \SI{3.49}{\hertz} ($\zeta_{\mathrm{t}} =
\SI{4.37}{\percent}$) and argued, on the basis of its generator-dominated mode
shape, that it should be excitable from the electrical port. This section
provides the experimental verification promised there and in
\S\ref{subsec:of-discussion}. The same electrical forcing experiment is run in
two models: first the reduced two-mass model, where the electromagnetic-torque
path is open by construction, and then the high-fidelity OpenFAST co-simulation,
where the mode is reproduced in the FMU drivetrain-interface layer. In both
cases a periodic electrical disturbance at the modal frequency is shown to drive
the shaft torque into a growing, sustained resonance, while an off-resonance
disturbance of equal amplitude does not.

\subsection{Forcing setup}
\label{subsec:fr-setup}

The electrical disturbance is applied as an oscillating shunt load at the
wind-turbine terminal bus. A shunt admittance of magnitude $y_{\mathrm{load}} =
\hat{P}/S_{\mathrm{base}}$ is modulated sinusoidally and switched on at a time
$t_{\mathrm{on}}$ (before which the system sits at its power-flow equilibrium),
\begin{equation}
  y_{\mathrm{sh}}(t) = y_{\mathrm{load}}\,
  \sin\!\big(2\pi f\,(t-t_{\mathrm{on}})\big)\,\mathbb{1}[t \ge t_{\mathrm{on}}].
  \label{eq:fr-forcing}
\end{equation}
The modulated load makes the wind-turbine terminal power oscillate,
$P_e(t) = P_0 + \hat{P}\sin(2\pi f\,t)$ as in \eqref{eq:Pe-forcing}, which
through $T_e = P_e/(\omega_{e,\mathrm{filt}}\eta)$ produces an oscillating
electromagnetic torque on the generator mass. The forcing frequency $f$ is swept
across the modal frequency, and for each run the resonance is quantified by the
\emph{post-onset peak-to-peak shaft torque}, $\Delta T_{\mathrm{sh}} =
\max_{t\ge t_{\mathrm{on}}} T_{\mathrm{sh}} - \min_{t\ge t_{\mathrm{on}}}
T_{\mathrm{sh}}$, complemented by the envelope build-up ratio and an FFT of the
response.

\subsection{Reduced two-mass model}
\label{subsec:fr-simplified}

With a forcing amplitude of $\hat{P} = \SI{2.0}{\mega\watt}$, the shaft-torque
response is dramatically frequency-selective. Off resonance (a
\SI{1.0}{\hertz} disturbance) the response is small and its envelope
\emph{decays} after onset (build-up ratio $0.39$), i.e.\ the disturbance merely
pushes the well-damped low-frequency dynamics. On resonance
(\SI{3.49}{\hertz}) the response is an order of magnitude larger and its
envelope is \emph{sustained} (build-up ratio $\approx 1.0$), the signature of a
forcing that continuously feeds the lightly damped mode
(Figure~\ref{fig:fr-onoff-simple}). An FFT of the on-resonance response peaks
exactly at the forcing frequency, confirming that it is the torsional mode, and
not a harmonic, that is being driven. Figure~\ref{fig:fr-buildup-simple} shows
the two-panel time history: the sinusoidal electrical forcing (top) and the
resonant build-up of the shaft torque towards its saturated amplitude (bottom).

\begin{figure}[t]
  \centering
  % \includegraphics[width=0.8\linewidth]{fig/WT_LEOGO_torsional_timeseries_on_off.png}
  \caption{Reduced model: shaft-torque response to an electrical disturbance of
  equal amplitude, on resonance (\SI{3.49}{\hertz}) versus off resonance
  (\SI{1.0}{\hertz}). The off-resonance response decays after onset while the
  on-resonance response is sustained.}
  \label{fig:fr-onoff-simple}
\end{figure}

\begin{figure}[t]
  \centering
  % \includegraphics[width=0.8\linewidth]{fig/WT_LEOGO_torsional_buildup.png}
  \caption{Reduced model: electrical forcing at \SI{3.49}{\hertz} (top, terminal
  power) and the resonant build-up of the shaft torque towards saturation
  (bottom).}
  \label{fig:fr-buildup-simple}
\end{figure}

Sweeping the forcing frequency traces out the resonance curve of
Figure~\ref{fig:fr-curve-simple}, with the values listed in
Table~\ref{tab:fr-sweep}. The peak-to-peak shaft torque rises from
\SI{0.22}{\mega\newton\metre} at \SI{1.0}{\hertz} to
\SI{7.96}{\mega\newton\metre} at \SI{3.49}{\hertz} --- an amplification of about
$36\times$ --- and falls off again above the mode, with the peak located
precisely at the eigenfrequency predicted by the modal analysis. The
sharpness of the curve is consistent with the quality factor $Q =
1/(2\zeta_{\mathrm{t}}) \approx 11$ of \eqref{eq:Q}.

\begin{figure}[t]
  \centering
  % \includegraphics[width=0.8\linewidth]{fig/WT_LEOGO_torsional_resonance_curve.png}
  \caption{Reduced model: peak-to-peak shaft torque versus electrical forcing
  frequency ($\hat{P}=\SI{2.0}{\mega\watt}$). The dashed line marks the
  torsional eigenfrequency; the response peaks there.}
  \label{fig:fr-curve-simple}
\end{figure}

\begin{table}[t]
  \centering
  \caption{Reduced-model frequency sweep: post-onset peak-to-peak shaft torque
  and envelope build-up ratio ($\hat{P}=\SI{2.0}{\mega\watt}$).}
  \label{tab:fr-sweep}
  \begin{tabular}{lcc}
    \toprule
    $f$ [\si{\hertz}] & $\Delta T_{\mathrm{sh}}$ [\si{\mega\newton\metre}] &
      Envelope ratio \\
    \midrule
    1.00 & 0.22 & 0.39 \\
    2.00 & 0.66 & 0.37 \\
    2.50 & 1.23 & 0.50 \\
    3.00 & 2.65 & 0.66 \\
    3.25 & 4.46 & 0.84 \\
    \textbf{3.49} & \textbf{7.96} & \textbf{1.01} \\
    3.75 & 5.53 & 0.87 \\
    4.00 & 4.09 & 0.76 \\
    5.00 & 2.35 & 0.75 \\
    \bottomrule
  \end{tabular}
\end{table}

\subsection{OpenFAST-FMU model}
\label{subsec:fr-fmu}

The experiment was then repeated in the high-fidelity co-simulation, where the
IEA \SI{15}{\mega\watt} turbine runs as an FMU coupled to the grid. A word on
where the torsional mode lives is needed to reconcile this with the negative
result of \S\ref{subsec:of-electrical}. Inside OpenFAST the drivetrain is kept
rigid (\texttt{DrTrDOF} disabled), and ROSCO owns the generator torque and
rejects external commands, so the OpenFAST \emph{structure} cannot be driven
electrically. The co-simulation interface layer, however, represents the
electrical machine as a lumped generator mass $H_e$ coupled through a compliant
shaft ($K_{pu},D_{pu}$) to the rotor speed \emph{prescribed} by the FMU
(\texttt{RotSpeed}, itself governed by ROSCO). The electromagnetic torque
$T_e = P_e/(\omega_{e,\mathrm{filt}}\eta)$ acts on this generator mass exactly
as in the reduced model, so the electrical$\to$mechanical path into the
\emph{interface} torsional mode is open by construction, independent of ROSCO,
which merely sets the rotor-side boundary.

Because the rotor speed is a prescribed boundary rather than a free inertia, the
interface reduces to a single mass on a spring, whose natural frequency is
\begin{equation}
  f_n = \frac{1}{2\pi}\sqrt{\frac{K_{pu}}{2H_e}} = \SI{3.485}{\hertz},
  \qquad
  \zeta = \frac{D_{pu}}{2\sqrt{2H_e K_{pu}}} = \SI{4.4}{\percent},
  \label{eq:fr-fmu-fn}
\end{equation}
using the running wrapper constants $H_e = \SI{0.0384}{\second}$,
$K_{pu} = 36.8$, $D_{pu} = 0.149$. This is the same \SI{3.49}{\hertz} torsional
mode, now reached from the single-mass limit of the two-mass system.

With a smaller forcing amplitude of $\hat{P} = \SI{0.5}{\mega\watt}$, the FMU
sweep (Table~\ref{tab:fr-fmu-sweep}, Figure~\ref{fig:fr-curve-fmu}) again shows
a clean resonance peaking at \SI{3.49}{\hertz}, where the peak-to-peak shaft
torque reaches \SI{8.57}{\mega\newton\metre}. The on- versus off-resonance time
series and the resonant build-up (Figures~\ref{fig:fr-onoff-fmu}
and~\ref{fig:fr-buildup-fmu}) reproduce the reduced-model behaviour inside the
full aeroelastic co-simulation. The demonstration therefore holds not only in
the idealised model but also in the detailed one, provided the excitation
targets the electrically coupled generator mass rather than the
controller-governed rotor torque.

\begin{table}[t]
  \centering
  \caption{OpenFAST-FMU frequency sweep: post-onset peak-to-peak shaft torque
  ($\hat{P}=\SI{0.5}{\mega\watt}$).}
  \label{tab:fr-fmu-sweep}
  \begin{tabular}{lc}
    \toprule
    $f$ [\si{\hertz}] & $\Delta T_{\mathrm{sh}}$ [\si{\mega\newton\metre}] \\
    \midrule
    1.00 & 1.65 \\
    2.00 & 1.79 \\
    2.50 & 1.98 \\
    3.00 & 3.51 \\
    3.25 & 5.54 \\
    \textbf{3.49} & \textbf{8.57} \\
    3.75 & 5.48 \\
    4.50 & 2.81 \\
    \bottomrule
  \end{tabular}
\end{table}

\begin{figure}[t]
  \centering
  % \includegraphics[width=0.8\linewidth]{fig/WT_LEOGO_FMU_torsional_resonance_curve.png}
  \caption{OpenFAST-FMU model: peak-to-peak shaft torque versus electrical
  forcing frequency ($\hat{P}=\SI{0.5}{\mega\watt}$). The response peaks at the
  \SI{3.49}{\hertz} torsional mode.}
  \label{fig:fr-curve-fmu}
\end{figure}

\begin{figure}[t]
  \centering
  % \includegraphics[width=0.8\linewidth]{fig/WT_LEOGO_FMU_torsional_timeseries_on_off.png}
  \caption{OpenFAST-FMU model: shaft torque on resonance (\SI{3.49}{\hertz})
  versus off resonance (\SI{1.0}{\hertz}) for an equal-amplitude electrical
  disturbance.}
  \label{fig:fr-onoff-fmu}
\end{figure}

\begin{figure}[t]
  \centering
  % \includegraphics[width=0.8\linewidth]{fig/WT_LEOGO_FMU_torsional_buildup.png}
  \caption{OpenFAST-FMU model: electrical forcing at \SI{3.49}{\hertz} (top) and
  the resonant build-up of the wrapper shaft torque (bottom).}
  \label{fig:fr-buildup-fmu}
\end{figure}

\subsection{Comparison of the two models}
\label{subsec:fr-comparison}

Figure~\ref{fig:fr-comparison} overlays the two resonance curves, each
normalised to its own peak (the absolute amplitudes are not directly comparable
because the models are driven at different forcing amplitudes ---
\SI{2.0}{} versus \SI{0.5}{\mega\watt} --- and have different inertia scalings).
The decisive point is that both curves peak at the \emph{same} frequency,
\SI{3.49}{\hertz}: the reduced model and the high-fidelity co-simulation locate
the identical torsional mode, and both respond to electrical forcing there.

Two physically meaningful differences remain. First, per unit of forcing power
the FMU response is much larger (\SI{8.57}{\mega\newton\metre} from
\SI{0.5}{\mega\watt}, versus \SI{7.96}{\mega\newton\metre} from
\SI{2.0}{\mega\watt}): in the single-mass limit \eqref{eq:fr-fmu-fn} the entire
compliant motion is concentrated in the small generator inertia $H_e$, so a
given torque disturbance produces a larger twist. Second, the FMU curve has a
higher off-resonance floor and thus appears less sharply selective. This is a
direct consequence of the prescribed rotor: with $\omega_m$ held fixed by the
FMU, a low-frequency electrical disturbance twists the shaft quasi-statically
against an immovable rotor, whereas the reduced model has a free rigid-body mode
that lets both masses drift together and relieves the low-frequency twist. The
resonant peak itself, and the mode location, are unaffected.

\begin{figure}[t]
  \centering
  % \includegraphics[width=0.8\linewidth]{fig/WT_LEOGO_torsional_model_comparison.png}
  \caption{Normalised resonance curves of the reduced two-mass model and the
  OpenFAST-FMU model, each scaled to its own peak. Both locate the same
  \SI{3.49}{\hertz} torsional mode; the FMU curve has a higher off-resonance
  floor because its rotor speed is a prescribed boundary.}
  \label{fig:fr-comparison}
\end{figure}

\subsection{Conclusion}
\label{subsec:fr-conclusion}

The forced-response experiments confirm the central hypothesis: a lightly
damped drivetrain torsional mode can be driven to a sustained mechanical
resonance from the electrical port. The reduced model demonstrates the effect
cleanly ($36\times$ amplification at \SI{3.49}{\hertz}, sustained envelope,
FFT peak at the forcing frequency), and the OpenFAST co-simulation reproduces
it at the same frequency once the excitation is applied to the electrically
coupled generator mass rather than to the ROSCO-governed rotor torque. It is
worth stressing that the \SI{3.49}{\hertz} mode is the compliant coupling of the
co-simulation interface, not OpenFAST's internal drivetrain torsion: enabling
OpenFAST's own \texttt{DrTrDOF} would introduce a different, much stiffer
($\sim\SI{30}{\hertz}$) mode lying outside the co-simulation bandwidth
(above the Nyquist frequency of the $\Delta t = \SI{0.01}{\second}$ step), and
is neither necessary nor helpful for the present demonstration.
